% Generated by report/make_restated.py from theory/B2_effective_curvature.tex: do not edit (edit the note and rerun).
\begin{restated}{Remark}{rem:B2-gd-rounding}{B2}{0}{rounding errors and the feature span: a heuristic}
This remark is a heuristic, not a proved statement. It concerns gradient descent in the situation of
Remark~\ref{rem:B2-gd-scope}(ii): logistic regression with $d>n$, $\theta_0\in S$ and $E_{\min}=S^\perp$. \emph{What is
proved} is the exact-arithmetic invariance of Lemma~\ref{lem:B2-span}(b), the alignment of Proposition~\ref{prop:B2-gd}(d)
for a quadratic, and the following bound. Write the computed iterates as
$\theta_t=\theta_{t-1}-\gamma\Fp(\theta_{t-1})+\delta_t$, where $\delta_t$ collects all the rounding errors of step $t$
(including those made in computing the gradient). Since $\Fp(\theta)_\perp=\lambda\theta_\perp$ for every $\theta$ and
$\etastar\in S$ (Lemma~\ref{lem:B2-span}(a)), the $S^\perp$-component of the error satisfies exactly
$v_{t,\perp}=(1-\gamma\lambda)v_{t-1,\perp}+\delta_{t,\perp}$; hence, for $\theta_0\in S$ and $0<\gamma\lambda\le1$ (true for
$\gamma\le1/L$),
\[
\norm{v_{t,\perp}}\;\le\;\sum_{s=1}^t(1-\gamma\lambda)^{t-s}\norm{\delta_s}\;\le\;\frac{\max_{s\le t}\norm{\delta_s}}{\gamma\lambda} .
\]
\emph{What is not proved} is the size of $\delta_t$. In the standard model of floating-point arithmetic it is of order
$\varepsilon_{\mathrm{mach}}\approx10^{-16}$ times the size of the quantities that are added ($\norm{\theta_{t-1}}$, and
$\gamma$ times the individual terms of the gradient), up to factors depending on $n$ and $d$; in particular it does not
decrease with the error $\norm{v_{t-1}}$. Under this model the component of the error along $E_{\min}$ stays below a level
of order $\varepsilon_{\mathrm{mach}}\max_s\norm{\theta_s}/(\gamma\lambda)$ (and is typically much smaller, since rounding
errors do not add up coherently), while the component in $S$ keeps decreasing geometrically. The former can therefore
become comparable with the error only when the error is within a factor of order $1/(\gamma\lambda)$ of the rounding
level; above that level one observes the behaviour inside $S$, that is, the alignment of
Proposition~\ref{prop:B2-gd}(d) with the bottom eigenspace of $H|_S$ in the quadratic model (for logistic regression, the
limit described, itself heuristically, in Remark~\ref{rem:B2-gd-heuristic}). Consistently, in the sanity check of
Remark~\ref{rem:B2-gd-heuristic} the $S^\perp$-component of the error stays at the rounding level $10^{-14}$. No proved
statement of this note depends on this remark; prediction (P1) below uses it only when $S\neq\R^d$ (e.g.\ $d>n$) and
$\theta_0\in S$, a situation that does not occur in the experiments of the project, where the features span $\R^d$
(Remark~\ref{rem:B2-P2-outcome}).
\end{restated}
